%% Section 5: Experiments
%% English manuscript section.
%% -----------------------------------------------------------------------

\section{Experiments}
\label{sec:experiments}

\subsection{Main Results}
\label{sec:main-results}

Table~\ref{tab:main} shows test-set performance for all model sizes and variants. IntSeqBERT
consistently outperforms both baselines on all scales and metrics. At the Large scale,
IntSeqBERT surpasses the Vanilla baseline by $+8.9$\,pt in Mag Acc and $+4.5$\,pt in MMA.
The Ablation model shows the largest MMA drop at Large scale ($-15.2$\,pt), directly
quantifying the contribution of the arithmetic-periodicity features. Notably, the Ablation
maintains non-trivial Mag Acc, but the remaining gap confirms that modulo information contributes
meaningfully to magnitude regression as well, while being essential for modulo prediction.

\begin{table}[htbp]
\caption{Test results. Mag Acc (\%), Sign Acc (\%), MMA (\%). \textbf{Bold} indicates the best value within each size group.}
\label{tab:main}
\centering
\begin{tabular}{llrrr}
\hline
Size   & Model          & Mag Acc        & Sign Acc       & MMA            \\
\hline
Small  & \textbf{IntSeq}& \textbf{94.73} & \textbf{97.78} & \textbf{40.43} \\
Small  & Vanilla        & 85.73          & 96.91          & 36.21          \\
Small  & Ablation       & 93.72          & 97.39          & 25.97          \\
Middle & \textbf{IntSeq}& \textbf{95.71} & \textbf{98.34} & \textbf{46.88} \\
Middle & Vanilla        & 87.37          & 97.42          & 42.53          \\
Middle & Ablation       & 92.45          & 97.90          & 31.93          \\
Large  & \textbf{IntSeq}& \textbf{95.85} & \textbf{98.54} & \textbf{50.38} \\
Large  & Vanilla        & 86.97          & 97.66          & 45.85          \\
Large  & Ablation       & 89.70          & 98.29          & 35.22          \\
\hline
\end{tabular}
\end{table}

\textbf{Learning curves.} Fig.~\ref{fig:learning_curve} shows validation loss across all
variants. Large IntSeqBERT converges from 2.17 to 1.01 over 200 epochs with no overfitting
(Train~1.00 / Val~1.01). Vanilla (1.77) and Ablation (1.39) maintain higher final loss.

\begin{figure}[htbp]
  \centering
  \includegraphics[width=\linewidth]{figures/fig4_learning_curves}
  \caption{Validation loss learning curves for all scales (Small / Middle / Large) and all
           variants. IntSeqBERT (solid blue) consistently achieves lower loss than Vanilla
           (dashed orange) and Ablation (dash-dot green). At the Large scale, IntSeqBERT
           converges to Val Loss = 1.01 at epoch 200.}
  \label{fig:learning_curve}
\end{figure}

\subsection{Magnitude Prediction}
\label{sec:magnitude}

\textbf{Scale-stratified analysis.} Table~\ref{tab:scale} reports per-bucket MSE on the test
set for Large models. The Vanilla model suffers catastrophic degradation in the Large bucket
(MSE = 2.10, $13\times$ that of IntSeqBERT), caused by all out-of-vocabulary integers being
absorbed into the \texttt{UNK} token. IntSeqBERT achieves the best MSE in all buckets except
Small, with particularly pronounced advantage at Medium and above. The Ablation degrades in the
Medium bucket without modulo context (MSE = 0.116 vs.\ IntSeqBERT 0.051), and collapses more
severely in the Huge and Astronomical buckets, indicating that FiLM modulation by the modulo
stream acts as an arithmetic structural constraint on magnitude estimation for large integers.

\begin{table}[htbp]
\caption{Scale-stratified MSE on the test set (Large models). Lower is better.}
\label{tab:scale}
\centering
\begin{tabular}{lrrr}
\hline
Bucket        & IntSeq         & Vanilla & Ablation \\
\hline
Small         & 0.111          & 0.138   & \textbf{0.103} \\
Medium        & \textbf{0.051} & 0.071   & 0.116   \\
Large         & \textbf{0.162} & 2.100   & 0.381   \\
Huge          & \textbf{2.082} & 22.73   & 5.021   \\
Astronomical  & \textbf{110.4} & 840.0   & 532.6   \\
\hline
\end{tabular}
\end{table}

Fig.~\ref{fig:scatter} shows predicted vs.\ true magnitude scatter plots for the three Large
model variants. IntSeqBERT achieves the highest coefficient of determination ($R^2 = 0.988$),
substantially reducing off-diagonal dispersion in the Large and Huge buckets compared to
Vanilla ($R^2 = 0.943$), visually corroborating the MSE results of Table~\ref{tab:scale}.

\begin{figure}[htbp]
  \centering
  \includegraphics[width=\linewidth]{figures/fig5_magnitude_scatter}
  \caption{Predicted vs.\ true magnitude ($\log_{10}$ scale) for Large models. Points are
           coloured by bucket (Small=blue circle, Medium=green circle, Large=yellow-orange
           square, Huge=red triangle, Astronomical=purple diamond). IntSeqBERT achieves
           $R^2 = 0.988$ vs.\ Vanilla $R^2 = 0.943$; Vanilla shows pronounced scatter above
           the Large bucket.}
  \label{fig:scatter}
\end{figure}

\subsection{Modulo Spectrum Analysis}
\label{sec:modulo-spectrum}

Fig.~\ref{fig:nig_spectrum} shows the NIG spectrum for $m = 2, \ldots, 101$ using Large
models. IntSeqBERT (solid blue) exceeds both Vanilla and Ablation across the entire spectrum,
with a visually clear contrast between prime and composite moduli.

\begin{figure}[h]
  \centering
  \includegraphics[width=\linewidth]{figures/fig2_nig_spectrum}
  \caption{NIG spectrum for moduli $m = 2, \ldots, 101$ (Large models). Grey shading marks
           prime moduli. The 95\% CI for IntSeqBERT (light blue band) is computed by
           bootstrapping.}
  \label{fig:nig_spectrum}
\end{figure}

\textbf{Finding 1: NIG is strongly negatively correlated with Euler's totient ratio.}
A Pearson correlation of $r = -0.851$ ($p < 10^{-28}$) is observed between NIG and
$\varphi(m)/m = \prod_{p \mid m}(1 - 1/p)$ (Fig.~\ref{fig:nig_phi}). Composite moduli with
many small prime factors achieve higher NIG, consistent with a CRT aggregation effect: when
$m$ is a common multiple of smaller moduli $m_1, m_2, \ldots$, then $x \bmod m$ encodes the
information of all those moduli simultaneously.
The highest NIG across all models and scales is achieved at $m = 96 = 2^5 \times 3$
($\varphi(96)/96 = 1/3$; Large IntSeq: NIG = 0.629, 95\% CI [0.622, 0.634]), consistent with
$m = 96$ being the largest modulus with totient ratio $1/3$ among $\{12, 24, 48, 72, 96\}$ and
thus distinguishing the widest value range. The exception is $m = 2$ (prime, NIG = 0.628,
second highest), reflecting the corpus-wide ubiquity of parity in OEIS sequences.

While this correlation is an empirical observation rather than a theorem, it carries a
design implication for the modulo stream itself: composite moduli reflect constraints from
several prime factors and prime powers at once, so the CRT aggregation effect lets a model
extract more number-theoretic information from fewer moduli. This gives a guideline for
choosing moduli---for example, when extending the spectrum beyond $m = 101$, adding
composite numbers with many small prime factors (such as highly composite numbers) is
likely to be more effective than simply adding more primes.

\begin{figure}[htbp]
  \centering
  \includegraphics[width=0.72\linewidth]{figures/fig2b_nig_vs_phi}
  \caption{NIG vs.\ Euler's totient ratio $\varphi(m)/m$ (Large IntSeqBERT). Composite
           moduli (blue circles, shade proportional to $m$) and prime moduli (red triangles)
           are shown separately. The regression line (grey dashed) indicates Pearson correlation of
           $r = -0.851$ ($p < 10^{-28}$). Notable moduli $m = 2$ (parity),
           $m = 60$ (Babylonian number), and $m = 96$ ($2^5 \times 3$ composite) are annotated.}
  \label{fig:nig_phi}
\end{figure}

\textbf{Finding 2: Parity (mod 2) accuracy stratifies models.}
Mod-2 accuracy at the Large scale is 85.65\% (IntSeq), 81.40\% (Vanilla), and 72.13\%
(Ablation), making parity the single modulus most sensitive to the modulo stream ($-13.5$\,pt
when removed). Table~\ref{tab:moduli} lists representative modulus accuracies.

\begin{table}[h]
\caption{Representative modulus accuracies for Large models (\%).}
\label{tab:moduli}
\centering
\begin{tabular}{rrrrl}
\hline
$m$ & IntSeq & Vanilla & Ablation & Interpretation \\
\hline
  2 & 85.65  & 81.40   & 72.13    & Parity \\
  3 & 72.62  & 65.22   & 53.72    & Ternary residue \\
  5 & 60.37  & 50.07   & 42.63    & Last decimal digit modulo 5 \\
 10 & 58.38  & 49.25   & 39.47    & Least significant decimal digit \\
 60 & 53.97  & 47.87   & 35.12    & Babylonian number (highly composite) \\
 96 & 51.82  & 47.29   & 34.44    & Composite ($2^5 \times 3$) \\
100 & 48.51  & 45.60   & 33.51    & Centesimal residue \\
\hline
\end{tabular}
\end{table}

\subsection{Solver-Based Next-Term Prediction}
\label{sec:solver-eval}

The Solver reconstructs candidate integers from the model's magnitude, sign, and modulo
predictions and ranks them by likelihood. We evaluate exact-match accuracy on 10,000 test
samples, where each sample corresponds to one OEIS sequence and the target is the \emph{last
term} (all preceding terms are given as context).

\emph{Valid-candidate rate} is the fraction of samples for which the Solver returns at least
one candidate. Vanilla always returns a candidate from its vocabulary softmax (rate = 100\%).
IntSeqBERT and Ablation may fail to find a valid candidate within the search range
(\texttt{none} mode; 13.4\% of Large-IntSeq calls), reducing their valid-candidate rate.

Table~\ref{tab:solver-eval} shows Solver evaluation results. Large-scale IntSeqBERT achieves
a Top-1 accuracy of 19.09\%, a \textbf{7.4-fold} improvement over the Vanilla
baseline (2.59\%), counting \texttt{none} returns as misses.

\begin{table}[htbp]
\caption{Solver evaluation: Top-1 and Top-10 exact-match accuracy (\%) and valid-candidate rate.}
\label{tab:solver-eval}
\centering
\begin{tabular}{llrrrr}
\hline
Size   & Model           & Top-1          & Top-10         & Sign Acc       & Valid rate (\%) \\
\hline
Small  & \textbf{IntSeq} & \textbf{14.05} & \textbf{21.00} & \textbf{98.73} & 90.59 \\
Small  & Vanilla         & 2.43           & 3.24           & 92.92          & 100.0 \\
Small  & Ablation        & 7.42           & 17.33          & 98.50          & 90.17 \\
Middle & \textbf{IntSeq} & \textbf{17.02} & \textbf{22.62} & \textbf{99.02} & 86.31 \\
Middle & Vanilla         & 2.43           & 3.41           & 92.71          & 100.0 \\
Middle & Ablation        & 9.88           & 20.52          & 98.74          & 90.34 \\
Large  & \textbf{IntSeq} & \textbf{19.09} & \textbf{26.23} & \textbf{99.02} & 86.64 \\
Large  & Vanilla         & 2.59           & 3.80           & 92.05          & 100.0 \\
Large  & Ablation        & 11.75          & 21.79          & 98.94          & 86.99 \\
\hline
\end{tabular}
\end{table}

\textbf{Accuracy by magnitude bucket.} Table~\ref{tab:solver-bucket} compares Solver
Top-1 accuracy across buckets for all three Large models. IntSeqBERT achieves 68.34\%
Top-1 in the Small bucket, while the Ablation reaches 54.55\% and Vanilla only 14.11\%.
Vanilla collapses to 0\% for Medium and above due to the UNK-token degradation of its
magnitude predictions, whereas IntSeqBERT retains meaningful accuracy up to the Medium
bucket (20.82\%).

\begin{table}[htbp]
\caption{Solver Top-1 / Top-10 accuracy (\%) by magnitude bucket (Large models).
         \textbf{Bold}: best per bucket.}
\label{tab:solver-bucket}
\centering
\begin{tabular}{l rr rr rr}
\hline
 & \multicolumn{2}{c}{\textbf{IntSeqBERT}} & \multicolumn{2}{c}{Vanilla} & \multicolumn{2}{c}{Ablation} \\
Bucket & Top-1 & Top-10 & Top-1 & Top-10 & Top-1 & Top-10 \\
\hline
Small        & \textbf{68.34} & \textbf{88.50} & 14.11 & 20.71 & 54.55 & 86.92 \\
Medium       & \textbf{20.82} & \textbf{31.50} &  0.00 &  0.00 &  5.61 & 18.88 \\
Large        & \textbf{ 0.31} & \textbf{ 0.67} &  0.00 &  0.00 &  0.03 &  0.05 \\
Huge         & \textbf{ 0.09} & \textbf{ 0.18} &  0.00 &  0.00 &  0.00 &  0.00 \\
Astronomical &  0.00          &  0.00          &  0.00 &  0.00 &  0.00 &  0.00 \\
\hline
\end{tabular}
\end{table}

\textbf{Mode breakdown (Large IntSeqBERT).} Dense mode (24.0\% of calls) achieves the highest
Top-1 accuracy at 61.06\%; Sieve mode (36.7\%) achieves 5.36\%; the zero-prediction shortcut
(2.7\%) achieves 89.96\%. CRT mode (23.2\%) achieves near-zero accuracy (0.09\%); 13.4\% of
calls return no valid candidate. The CRT limitation is discussed in Section~\ref{sec:limitations}.
